\documentclass[11pt,
aspectratio=23, % для экранов 2:3
]{beamer} 

\setbeamersize{sidebar width right=1cm}

\input{vkr_identification}  % персональные данные
\usepackage{tex/PnuPresentStyle}     % не трогать 

\begin{document}

\begin{titleframe}
\titlepage
\end{titleframe}

\section{Введение}
\begin{frame}{Постановка задачи}
Какой-то текст и ссылка на литературу \cite{dyakonov}.
\end{frame}

\section{Методика}
\begin{frame}{Основные формулы}
 Вот эта формула:
   \begin{equation}
       \dfrac{D\underline{u}}{Dt} = \dfrac{1}{\rho} \underline{\nabla} \cdot \left( -p\underline{I} + \underline{\tau} \right) + \underline{f}.
   \end{equation}
\end{frame}

\begin{frame}{Слайд с картинкой}
 
Пример рисунка и ссылки на рисунок \ref{fig:figure1} на слайде \insertframenumber.

\begin{figure}[ht!] %[h!]
	\begin{center}
		\includegraphics[scale=0.2]{fig_template}
	\end{center}
	\caption{Пример оформления}
\label{fig:figure1}
\end{figure}
\end{frame}

\section{Результаты}
\begin{frame}{Интерпретация результатов}
  Результаты опубликованы тут \cite{pakin}.
\end{frame}

\section{Заключение}
\begin{frame}{Заключение}
Пример списка
    \begin{itemize}
        \item Раз
        \item Два
        \item Три
    \end{itemize}

\end{frame}

\section*{} % Not belong to any section
% REFERENCES
\begin{frame}
%	[allowframebreaks]
\frametitle{Библиография}
% \usebeamercolor{normal text}
\printbibliography[title={Библиография},heading=bibintoc]
\end{frame}

%%%%%%%%%% Благодарности %%%%%%%%%%%%%%

\begin{nobarframe}
\vfill
\begin{center}
\Large
Спасибо кому-нибудь .    
\end{center}
\vfill
\end{nobarframe}

\end{document}
